\chapter{总结与展望}
\label{chap:conclusion}

\section{研究成果总结}
\label{sec:achievement}

本文围绕“雨天环境下车道线检测鲁棒性不足”这一问题，按照“问题分析 - 数据构建 - 方法设计 - 工程验证”的思路开展研究，完成了从数据准备到边缘端部署验证的一体化工作。论文的主要研究成果可概括为以下四个方面。

\subsection{完成了雨天干扰机理与任务难点的系统梳理}

本文从雨滴遮挡、水膜反射和低对比度退化三个角度分析了雨天场景中车道线特征退化的主要来源，指出车道线检测在雨天条件下容易出现结构中断、伪边缘误激活和局部定位不稳定等问题。这一分析为后续数据集构建和模块设计提供了较明确的问题背景，也使全文研究始终围绕“提升雨天条件下车道线检测的可靠性与可部署性”展开。

\subsection{完成了 RainLane 数据集的整理、增强与实验组织}

本文构建并整理了自建 RainLane 数据集，并按照 CLRNet 可直接接入的方式完成了数据组织。当前确认的原始有效样本共 6017 张，其中训练集 4208 张、验证集 904 张、测试集 905 张。围绕雨纹干扰、水膜反射、局部遮挡和成像退化等典型问题，本文进一步设计了针对性的增强流程，将样本规模扩展至 24068 张，为第四章模型训练和第五章分场景评估提供了统一的数据基础。

\subsection{完成了基于 CLRNet 的轻量化抗干扰模块设计与验证}

在方法层面，本文以 CLRNet 为基线，设计了由 Frequency Gate 与 Directional Attention 组成的轻量化抗干扰方案，并完成了相应的实现与消融验证。实验结果表明，在 RainLane 测试集上，基线 CLRNet 的 F1 为 59.56\%，Frequency Gate 为 59.65\%，完整 FGM 为 59.70\%。这一结果说明，所提出的改进模块在不引入复杂外部预处理网络的前提下，能够对雨天干扰带来稳定收益；其中，完整 FGM 在 Precision 与 Recall 之间表现更为均衡，并在中雨等复杂场景中体现出一定优势。

\subsection{完成了 Orin 平台上的部署链路与性能验证}

在工程验证层面，本文完成了从 PyTorch 到 ONNX 再到 TensorRT 的部署链路，并在 Jetson Orin 平台上对已完成全链路导出的代表增强模型进行了测试。结果表明，PyTorch FP32、ONNX Runtime FP32、TensorRT FP16 和 TensorRT INT8 四种后端均已成功运行；其中 TensorRT FP16 的平均推理时延为 7.5435 ms、吞吐为 132.5641 FPS、峰值显存占用为 29.3867 MB，TensorRT INT8 的平均推理时延进一步降至 5.9027 ms、吞吐达到 169.4154 FPS，但其输出一致性弱于 FP16。结合第五章的定量与定性分析，当前更适合作为实际部署选择的是 TensorRT FP16 后端。

\section{创新点总结}
\label{sec:innovation}

基于本文的研究内容，主要创新点体现在以下三个方面。

\subsection{围绕雨天场景构建了可直接服务训练与评估的数据基础}

与仅少量包含雨天样本的通用车道线数据集不同，本文整理的 RainLane 直接面向“雨天车道线检测”这一研究目标，并在数据组织上与 CLRNet 训练流程保持兼容。通过对原始样本、增强样本和五类 RainLane 场景评测子集的统一整理，本文为离线精度评估、分场景分析与部署验证提供了同一套数据基础，使后续实验能够围绕同一问题展开。

\subsection{提出了面向雨天退化模式的轻量化抗干扰改进思路}

本文并未引入额外的大型去雨网络，而是从特征表达层面出发，设计了频域门控与方向注意力相结合的轻量化改进模块。该思路一方面针对雨纹噪声、反光伪边缘和细长结构断裂等雨天典型退化模式进行抑制，另一方面尽量保持原有检测框架的简洁性和部署友好性。实验结果表明，这一设计能够在较小结构改动下带来稳定的精度收益。

\subsection{形成了“数据 - 方法 - 部署”相贯通的研究路径}

本文的工作并非停留在单一模型改进，而是形成了从雨天问题分析、数据构建、模块设计到 Orin 部署验证的完整研究路径。第三章主要解决数据如何准备的问题，第四章围绕模块设计及其验证展开，第五章进一步给出部署过程与工程结果。按照这一结构展开，论文从研究动机到实验结论之间保持了较清晰的逻辑衔接。

\section{未来工作展望}
\label{sec:future_work}

尽管本文已经在数据、方法和部署三个层面完成了较系统的工作，但围绕模型性能、泛化能力与边缘端适配性仍有进一步深入的空间。未来工作可重点围绕以下几个方向展开。

\subsection{继续完善完整 FGM 的联合训练与系统验证}

后续工作可围绕完整 FGM 的联合训练策略继续展开，例如进一步优化残差缩放、分支权重与训练周期设置，并补充更系统的验证与部署评估，从而更充分地考察双分支结构在复杂雨天场景下的潜在收益。

\subsection{补充更完整的泛化与跨场景评估}

本文现阶段的定量分析主要基于 RainLane 数据集及其五类场景评测划分。后续可进一步补充跨数据集测试、多天气条件测试和更多极端干扰场景评估，以验证模型在不同域偏移条件下的稳定性，并更系统地分析频域抑噪与方向建模机制对不同退化模式的适应差异。

\subsection{推进边缘端部署的定量化与系统化测试}

第五章已经完成了 PyTorch、ONNX Runtime 与 TensorRT 多后端对比，并给出了时延、吞吐、显存与输出一致性结果。未来可在此基础上继续补充功耗、热稳定性、长时间运行抖动和视频流连续推理性能等指标，使部署结论从“能够运行”进一步提升到“能够长期稳定运行”的工程层面。

\subsection{扩展到更复杂天气与闭环应用场景}

本文当前重点聚焦于雨天场景，但真实 ADAS 系统仍需面向夜间强反光、雾天、雪天以及多天气叠加工况。后续研究可在现有工作基础上继续扩展多天气样本与统一建模框架，并尝试将车道线检测结果接入更完整的感知与控制链路中，评估其对下游车道保持或路径规划任务的影响，从而进一步提升研究结论的工程意义。

综上，本文围绕雨天车道线检测问题形成了较清晰的研究主线，并结合实验结果与 Orin 部署结果对主要结论进行了验证。后续工作将在此基础上继续完善双分支训练、泛化评估和系统级部署测试，使研究结论更加充分和稳健。
